\documentclass[10pt,twocolumn]{article}
\usepackage[T1]{fontenc}
\usepackage[utf8]{inputenc}
\usepackage{lmodern}
\usepackage[margin=1.8cm,top=2.4cm,bottom=2cm,columnsep=0.6cm]{geometry}
\usepackage{fancyhdr}
\usepackage{titlesec}
\usepackage{booktabs}
\usepackage{amsmath,amssymb,amsfonts}
\usepackage{microtype}
\usepackage{xcolor}
\usepackage{mdframed}
\usepackage{parskip}
\usepackage{enumitem}
\usepackage{hyperref}

\definecolor{rulered}{RGB}{180,30,30}
\definecolor{boxgray}{RGB}{245,245,245}
\definecolor{keywordgray}{RGB}{100,100,100}

\pagestyle{fancy}
\fancyhf{}
\fancyhead[L]{\textsc{\small Claude Mag}}
\fancyhead[C]{\small\textit{Machine Learning Research Digest}}
\fancyhead[R]{\small 2026-07-20}
\fancyfoot[C]{\small\thepage}
\renewcommand{\headrulewidth}{0.8pt}

\titleformat{\section}{\normalsize\bfseries\color{rulered}}{}{0pt}{}%
  [\vspace{-4pt}\color{rulered}\rule{\columnwidth}{0.4pt}]
\titlespacing{\section}{0pt}{8pt}{4pt}

\hypersetup{colorlinks=true,urlcolor=rulered,linkcolor=black}

\begin{document}
\setlength{\parindent}{0pt}
\setlength{\parskip}{4pt}

{\small\color{keywordgray}\textbf{Keywords:}
  reinforcement learning \textbullet{}
  zero RL \textbullet{}
  emergent reasoning \textbullet{}
  trillion-parameter scaling}\par\vspace{4pt}

{\LARGE\bfseries\raggedright
  Ring-Zero}\par\vspace{4pt}

{\large\itshape\raggedright
  Ant Group Scales Zero RL to One Trillion Parameters,
  Observing Five Emergent Reasoning Behaviors and 84.2\%
  on AIME 2026 Without Human-Annotated Data}\par\vspace{6pt}

{\small\textbf{By} Xinyu Tang, Qianggang Cao, Yurou Liu, Yuliang Zhan,
  Xiaochong Lan, Yifan Li, Yuchen Yan, Han Peng, Zican Dong,
  Zhenduo Zhang, et al.\ (Ant Group InclusionAI)}\par\vspace{2pt}

{\small\color{keywordgray}arXiv:2607.12395 \textbullet{}
  July 14, 2026}
\par\vspace{2pt}\noindent{\color{rulered}\rule{\columnwidth}{0.6pt}}\vspace{6pt}

Zero reinforcement learning --- RL with verifiable rewards and no
human-annotated data --- has demonstrated striking results at the 70B
scale, but whether its training dynamics and emergent capabilities
extend to far larger models remained unknown.
Ant Group's InclusionAI now closes that gap: Ring-Zero applies zero
RL to Ring-2.5-1T, a one-trillion-parameter hybrid linear-attention
model, and documents not only successful scaling but five qualitatively
new behaviors that appear only at this scale.

\begin{mdframed}[backgroundcolor=boxgray,linecolor=rulered,linewidth=1pt,
    innerleftmargin=6pt,innerrightmargin=6pt,
    innertopmargin=6pt,innerbottommargin=6pt]
\textbf{\small AT A GLANCE}\par\vspace{4pt}
\begin{description}[leftmargin=0pt,labelsep=4pt,itemsep=2pt,parsep=0pt,
    font=\small\bfseries]
  \item[Problem:] Naive scaling of zero RL collapses beyond $\sim$70B
    parameters due to readability degradation, token redundancy, and
    distribution drift.
  \item[Why it matters:] Trillion-parameter reasoning models are now
    practical; understanding and controlling zero RL at that scale is
    necessary for continued progress.
  \item[Method:] Clipped importance sampling, training-inference ratio
    correction, and mixed-precision control applied to Ring-2.5-1T (1T
    parameters, 1M context window).
  \item[Result:] 84.2\% on AIME 2026 with first-stage training only;
    five new emergent behaviors including self-verification and context
    anxiety.
  \item[Evidence:] Evaluated on AIME 2024/2025/2026, HMMT, and
    IMOAnswerBench; seven-benchmark average 74.8\%.
\end{description}
\end{mdframed}

\section{The Big Picture}

Zero RL uses only the binary signal of whether the model's answer
is correct to teach chain-of-thought reasoning --- no annotated
thought chains, no human preference data, no teacher demonstrations.
DeepSeek-R1-Zero showed that this surprisingly minimal training
signal is sufficient to produce sophisticated mathematical reasoning
at the 70B scale, but every subsequent study has operated at or below
that size.

The trillion-parameter regime matters because it is where
qualitatively new capabilities have historically emerged: GPT-4-class
reasoning, multimodal integration, and in-context meta-learning all
appeared near or above the 1T parameter threshold. Ring-Zero asks
whether the \emph{training dynamics} of zero RL also exhibit threshold
behavior at this scale --- and finds that they do.

\section{Why Existing Approaches Fall Short}

Three failure modes obstruct naive trillion-parameter zero RL.
First, \textbf{poor readability}: at large scale, unconstrained
policy gradient updates cause the model to generate verbose,
repetitive internal monologues that consume tokens without reasoning.
Second, \textbf{token redundancy}: the model learns to pad reasoning
traces with low-effort repetition, inflating apparent solution length
while contributing no additional correctness.
Third, \textbf{distribution drift}: at 1T parameters, rollout
distributions shift rapidly relative to the behavior policy used for
importance-weighted gradient estimation, destabilizing training after
a few hundred steps.

Existing stabilization techniques developed for 70B models
(gradient clipping, KL penalties) are insufficient at 1T because the
optimizer's effective learning rate --- after importance-weighted
corrections --- can grow dramatically when policy entropy is high in
a trillion-parameter space.

\section{Core Mechanism}

Ring-Zero addresses the three failure modes with three targeted
interventions. \textbf{Clipped importance sampling} restricts the
off-policy correction ratio to $[1-\epsilon, 1+\epsilon]$ per token,
preventing any single high-variance rollout from dominating the
gradient update. \textbf{Training-inference ratio correction}
compensates for the mismatch between the truncated context windows
used during RL rollouts and the full context available at inference.
\textbf{Mixed-precision control} uses BF16 for activations and FP32
for gradient accumulation, with selective FP8 for embedding layers
to keep the 1T model within GPU memory constraints.

\section{Technical Formulation}

The training-inference ratio correction scales the effective reward
by the ratio of inference to training context lengths:
\[
\hat{r}_t = r_t \cdot \left(\frac{C_{\text{inf}}}{C_{\text{train}}}\right)^\gamma
\]
where $C_{\text{inf}}$ is the inference context length, $C_{\text{train}}$
is the rollout context length, and $\gamma$ is a hyperparameter tuned
per model size (set to 0.4 for Ring-2.5-1T).

Clipped importance sampling modifies the standard importance weight:
\[
\tilde{\rho}_t = \min\!\left(\frac{\pi_\theta(a_t \mid s_t)}{\pi_{\theta_{\text{old}}}(a_t \mid s_t)},\; 1 + \epsilon\right)
\]
without a lower clip (since the lower clip would suppress on-policy
updates when the policy has improved).

\section{Learning Procedure}

Training reliably unfolds in two sequential phases identified across
multiple random seeds:

\begin{enumerate}[itemsep=2pt,parsep=0pt]
  \item \textbf{Discovery phase} (first $\approx$30\% of steps):
    The model actively expands reasoning boundaries, activating
    dormant pathway combinations that were not used in pre-training.
    Output diversity is high; format is unstable.
  \item \textbf{Sharpening phase} (remaining $\approx$70\% of steps):
    The model refines its policy within the stylistic and
    procedural boundaries established in discovery. Output quality
    and format consistency improve monotonically.
\end{enumerate}

The phase transition is sharp and reproducible: gradient entropy
drops by approximately 40\% between steps 800 and 1200 across all
runs, marking the boundary between discovery and sharpening.

\section{Five Emergent Behaviors}

Five behaviors appear at trillion-parameter scale that are absent at
70B, identified through systematic activation probing and output
pattern analysis:

\begin{enumerate}[itemsep=2pt,parsep=0pt]
  \item \textbf{Self-verification:} The model spontaneously checks
    intermediate conclusions before committing to a final answer,
    inserting verification sub-problems without being prompted to.
  \item \textbf{Context anxiety:} The model models itself as a
    resource-constrained agent operating within a finite context
    window, making strategic decisions about what to include based
    on this self-model.
  \item \textbf{Anthropomorphism:} The model adopts a persistent
    first-person problem-solving persona and maintains it across
    long reasoning chains.
  \item \textbf{Structured formatting:} Headers, numbered steps,
    and boxed answers emerge without any formatting reward.
  \item \textbf{Parallel reasoning:} The model explores multiple
    solution strategies in parallel within a single forward pass,
    then selects among them.
\end{enumerate}

\section{Experimental Findings}

\begin{center}
\small
\begin{tabular}{lcc}
\toprule
Benchmark & Ring-2.5-1T-Zero & Prior SOTA \\
\midrule
AIME 2026 & \textbf{84.2\%} & 78.5\% \\
AIME 2025 & 81.7\% & 79.3\% \\
AIME 2024 & 83.4\% & 82.1\% \\
HMMT & 78.3\% & 74.0\% \\
IMOAnswerBench & 61.2\% & 55.7\% \\
\midrule
7-Benchmark Avg. & \textbf{74.8\%} & 70.3\% \\
\bottomrule
\end{tabular}
\end{center}

Results are vendor-reported by InclusionAI; independent third-party
evaluation was not yet available at press time.

\section{Ablations and Interpretation}

\textbf{Without clipped importance sampling:} Training diverges
after approximately 500 steps at 1T scale; the effect is absent at
70B, confirming it is a trillion-parameter-specific failure mode.

\textbf{Without training-inference ratio correction:} AIME 2026 accuracy
drops from 84.2\% to 79.1\%, a 5.1 percentage point loss attributable
to underestimating long-context performance during training.

\textbf{Context anxiety phase transition:} The behavior is absent
below $\sim$800B parameters and strengthens above 1T, suggesting a
phase transition rather than a smooth scaling curve.

\textbf{Self-verification contribution:} Suppressing self-verification
tokens at inference reduces AIME 2026 accuracy by approximately 2--3
percentage points, confirming its functional role in correctness.

\section*{Reference}

Xinyu Tang et al.\
\textbf{Ring-Zero: Scaling Zero RL to a Trillion Parameters for
Emergent Reasoning}.
arXiv:2607.12395, July 2026.
\url{https://arxiv.org/abs/2607.12395}

\end{document}
